\section{Quản lý dự án và Phân công nhiệm vụ}

\subsection{Quy trình làm việc và Công cụ quản lý}

Để hiện thực hóa một hệ thống phức tạp từ tầng thu thập AI đến tầng lưu trữ không gian PostGIS, nhóm đã thiết kế một quy trình làm việc khoa học, tận dụng tối đa các công cụ cộng tác hiện đại:

\begin{itemize}
    \item \textbf{Quản lý tiến độ (Project Management):} Nhóm sử dụng nền tảng \textbf{Trello} để phân phối và giám sát công việc theo mô hình Kanban. Các tác vụ được chia nhỏ thành các thẻ (cards) với thời hạn cụ thể, giúp các thành viên bám sát tiến độ của từng giai đoạn: Phân tích, Thiết kế Schema, Hiện thực ETL và Đánh giá.
    \item \textbf{Quản lý mã nguồn và Phiên bản (Source Control):} Toàn bộ mã nguồn Python cho pipeline ETL, các module AI (YOLOv8x) và script khởi tạo cơ sở dữ liệu được quản lý tập trung trên \textbf{GitHub}. Nhóm tuân thủ quy trình Branching chặt chẽ để đảm bảo tính ổn định của mã nguồn khi có nhiều thành viên cùng tham gia hiện thực.
    \item \textbf{Thiết kế và Công cụ cơ sở dữ liệu:} Sử dụng \textbf{LucidChart} để phác thảo và chuẩn hóa mô hình Galaxy Schema. Khâu quản trị và truy vấn dữ liệu không gian được thực hiện thông qua \textbf{pgAdmin 4} và \textbf{DBeaver}, giúp tối ưu hóa việc kiểm tra tính đúng đắn của các bảng Dimension và Fact.
    \item \textbf{Cơ chế trao đổi và Báo cáo:}
    \begin{itemize}
        \item \textbf{Họp nội bộ (Daily Stand-up):} Nhóm duy trì các buổi họp hằng ngày (Online/Offline) để cập nhật tiến độ và giải quyết các điểm nghẽn kỹ thuật ngay lập tức.
        \item \textbf{Họp hướng dẫn (Weekly Meeting):} Nhóm thực hiện báo cáo định kỳ hằng tuần với GVHD (ThS. Bùi Tiến Đức) qua \textbf{Google Meet} để nhận góp ý chuyên môn và điều chỉnh Schema theo yêu cầu thực tế.
    \end{itemize}
\end{itemize}

\subsection{Chi tiết phân công công việc}

Bảng dưới đây tóm tắt vai trò chủ chốt của các thành viên trong nhóm:

% Thêm lệnh này vào trước bảng hoặc trong Preamble để X là căn giữa dọc
\renewcommand{\tabularxcolumn}[1]{m{#1}}

\begin{table}[htbp]
\centering
\renewcommand{\arraystretch}{1.5}
\begin{tabularx}{\textwidth}{|>{\centering\arraybackslash\bfseries}m{2.5cm}|>{\centering\arraybackslash}m{2cm}|X|>{\centering\arraybackslash}m{2cm}|}
\hline
\rowcolor[HTML]{EFEFEF} 
Thành viên & MSSV & \centering \textbf{Nhiệm vụ trọng tâm} & Hoàn thành \\ \hline

Nguyễn Thành Dũng & 2210589 & 
\begin{singlespace}
\begin{itemize}[leftmargin=*, nosep, before=\vspace{0.5pt}, after=\vspace{0pt}]
    \item Thiết kế Galaxy Schema sơ khai và hiện thực kiến trúc ETL chính.
    \item Cấu hình PostgreSQL/PostGIS và quản lý pipeline dữ liệu đa nguồn.
    \item Nghiên cứu xu hướng DW hiện đại để hoàn thiện Schema cuối cùng.
    \item Hiện thực Demo App.
\end{itemize} 
\end{singlespace} & 100\% \\ \hline

Hồ Trần Ngọc Liêm & 2211835 & 
\begin{singlespace}
\begin{itemize}[leftmargin=*, nosep, before=\vspace{0.5pt}, after=\vspace{0pt}]
    \item Phân tích Database VOH và cải tiến các bảng Dimension (Dim).
    \item Nghiên cứu và kết nối các nguồn API: TomTom, OSM, OWM.
    \item Đề xuất và tinh chỉnh hệ thống thuộc tính thuộc tầng không gian.
\end{itemize} 
\end{singlespace} & 100\% \\ \hline

Đinh Nguyễn Việt Khiêm & 2211565 & 
\begin{singlespace}
\begin{itemize}[leftmargin=*, nosep, before=\vspace{0.5pt}, after=\vspace{0pt}]
    \item Nghiên cứu tối ưu hóa Schema và tìm kiếm nguồn dữ liệu thứ cấp.
    \item Thiết kế và cải tiến bảng Fact dựa trên nghiệp vụ phân tích lưu lượng.
    \item Kiểm tra độ mịn dữ liệu (Granularity) và tính khả thi của thuộc tính.
\end{itemize} 
\end{singlespace} & 100\% \\ \hline

\end{tabularx}
\vspace{6mm}
\caption{Bảng phân công nhiệm vụ chi tiết và mức độ hoàn thành}
\end{table}

\subsection{Phân tích chi tiết đóng góp của các thành viên}

Dựa trên bảng phân chia tổng quát, mỗi thành viên đã tham gia sâu vào các khâu kỹ thuật cốt lõi của đồ án để đảm bảo tính nhất quán từ dữ liệu thô đến thông tin phân tích:

\subsubsection{Thành viên Nguyễn Thành Dũng: Kiến trúc sư ETL và Schema}

Với vai trò là kiến trúc sư chính, Nguyễn Thành Dũng chịu trách nhiệm định hình toàn bộ khung kỹ thuật và logic nghiệp vụ cho kho dữ liệu. Dũng đóng vai trò cầu nối giữa yêu cầu phân tích giao thông thực tế và việc hiện thực hóa các cấu trúc dữ liệu phức tạp trên nền tảng PostGIS cũng như hiện thực phần Demo App Saigon Traffic Monitor.

\begin{itemize}
    \item \textbf{Phân tích nghiệp vụ và Thiết kế Schema tiến hóa:} 
    \begin{itemize}[leftmargin=*, nosep]
        \item Chủ trì khâu phân tích luồng nghiệp vụ giao thông đô thị để xác định các thực thể chính, từ đó thiết kế bản thảo Schema sơ khai tập trung vào tính toàn vẹn dữ liệu.
        \item Nghiên cứu và áp dụng các xu hướng thiết kế Kho dữ liệu hiện đại để chuyển đổi từ Star Schema truyền thống sang mô hình Galaxy Schema kết hợp Snowflake Dimensions cho các chiều không gian.
        \item Tổng hợp kết quả từ khâu khảo sát nguồn dữ liệu và logic ETL để tinh chỉnh các thuộc tính, đảm bảo Schema cuối cùng tối ưu hóa được cả khâu lưu trữ lẫn khả năng mở rộng cho các bài toán dự báo trong tương lai.
    \end{itemize}

    \item \textbf{Hiện thực hóa lõi hệ thống ETL (Core ETL Implementation):} 
    \begin{itemize}[leftmargin=*, nosep]
        \item Xây dựng các script Python chủ đạo để điều phối quy trình trích xuất và biến đổi dữ liệu từ đa nguồn (AI, API, GIS).
        \item Thiết kế và triển khai cơ chế \textbf{Batch Load}, cho phép hệ thống nạp đồng thời hàng triệu bản ghi lưu lượng vào cơ sở dữ liệu mà không gây ra hiện tượng nghẽn cổ chai hay xung đột tài nguyên.
        \item Hiện thực các hàm xử lý logic nghiệp vụ phức tạp ngay trong quá trình nạp (như tính toán PCU, nhãn ca làm việc) để giảm tải cho tầng truy vấn báo cáo.
    \end{itemize}

    \item \textbf{Kỹ thuật Cơ sở dữ liệu và Tối ưu hóa hiệu năng (DB Engineering):} 
    \begin{itemize}[leftmargin=*, nosep]
        \item Trực tiếp khởi tạo và quản trị hệ thống cơ sở dữ liệu PostgreSQL tích hợp PostGIS, thiết lập các ràng buộc không gian chặt chẽ giữa các bảng.
        \item Cấu hình kỹ thuật \textbf{Partitioning} (phân vùng bảng Fact theo tháng), giúp duy trì hiệu năng truy vấn ở mức cao ngay cả khi dữ liệu tích lũy qua nhiều năm.
        \item Thiết lập hệ thống chỉ mục đa tầng, kết hợp giữa B-Tree cho dữ liệu định danh và \textbf{GIST Index} cho dữ liệu tọa độ không gian, đảm bảo các phép toán tìm kiếm lân cận hoặc Map Matching diễn ra với độ trễ tối thiểu.
    \end{itemize}
\end{itemize}

\paragraph{Ý nghĩa đóng góp:} Những nỗ lực của Dũng trong khâu thiết kế và hiện thực lõi đã tạo ra một nền tảng vững chắc, giúp hệ thống không chỉ vận hành ổn định trong thời gian thực mà còn cung cấp cấu trúc dữ liệu "sạch" và tối ưu cho các thành viên khác thực hiện các phép phân tích chuyên sâu.

\subsubsection{Thành viên Hồ Trần Ngọc Liêm: Chuyên gia Dữ liệu không gian và API}

Với vai trò là chuyên gia phụ trách tầng dữ liệu không gian và kết nối dịch vụ bên thứ ba, Hồ Trần Ngọc Liêm đóng đóng góp then chốt trong việc làm giàu ngữ cảnh và đảm bảo tính chính xác về mặt địa lý cho kho dữ liệu. Liêm đã trực tiếp xử lý các thách thức về sự khác biệt cấu trúc giữa các nguồn dữ liệu để tạo ra một hệ quy chiếu không gian thống nhất.

\begin{itemize}
    \item \textbf{Tối ưu hóa Dimension không gian và Tiền xử lý GIS:} 
    \begin{itemize}[leftmargin=*, nosep]
        \item Đi sâu vào phân tích cấu trúc nội hàm của các bảng \texttt{dim\_segment} và \texttt{dim\_location}, xác định các thuộc tính trọng yếu cần thiết cho việc mô hình hóa mạng lưới đường bộ phức tạp.
        \item Trực tiếp xử lý kỹ thuật tính toán góc \textbf{Azimuth} (hướng di chuyển) và chuẩn hóa các hệ tọa độ cho từng phân đoạn đường. Đây là bước tiền đề bắt buộc để thuật toán \textbf{Map Matching} có thể khớp chính xác dữ liệu vận tốc từ API vào đúng làn đường và hướng di chuyển thực tế.
        \item Phân tích và số hóa ranh giới hành chính đa tầng, hỗ trợ thiết lập cơ chế \textbf{Spatial Join} tự động để gán thông tin Phường/Quận cho từng tọa độ lưu lượng.
    \end{itemize}

    \item \textbf{Phát triển module kết nối API đa nguồn (Data Ingestion):} 
    \begin{itemize}[leftmargin=*, nosep]
        \item Nghiên cứu và hiện thực các module giao tiếp với \textbf{TomTom Traffic API} để trích xuất vận tốc di chuyển thực tế và vận tốc tự do (free-flow speed) theo thời gian thực.
        \item Xây dựng luồng thu thập dữ liệu từ \textbf{OpenWeatherMap (OWM)}, lấy các chỉ số về nhiệt độ, lượng mưa và tầm nhìn để tạo ra các biến giải thích (explanatory variables) cho bài toán phân tích tác động ngoại cảnh.
        \item Khai thác \textbf{OpenStreetMap (OSM)} để xây dựng hệ thống hạ tầng tĩnh, đảm bảo dữ liệu nền tảng luôn có độ bao phủ rộng và cập nhật.
    \end{itemize}

    \item \textbf{Phân tích Database nguồn và Tinh chỉnh thuộc tính Schema:} 
    \begin{itemize}[leftmargin=*, nosep]
        \item Trực tiếp khảo sát và phân tích cấu trúc dữ liệu từ nguồn VOH để tìm kiếm các đặc tính giao thông tiềm năng có thể tích hợp vào kho dữ liệu.
        \item Đề xuất và thực hiện cải tiến hệ thống thuộc tính cho các bảng Dimension (phiên bản 2), bổ sung các trường dữ liệu đặc trưng về loại đường và điều kiện hạ tầng dựa trên kết quả nghiên cứu các xu hướng DW-GIS hiện đại.
        \item Thiết lập mối liên kết logic giữa các chỉ số thời tiết và tốc độ di chuyển, hỗ trợ đắc lực cho việc xây dựng các bảng Fact có độ giàu thông tin cao.
    \end{itemize}
\end{itemize}

\paragraph{Ý nghĩa đóng góp:} Những nỗ lực của Liêm trong việc xử lý dữ liệu không gian và kết nối API đã giúp kho dữ liệu không chỉ dừng lại ở các con số thống kê khô khan mà trở thành một hệ thống giàu ngữ cảnh. Khả năng làm sạch và đối khớp dữ liệu GIS của Liêm là yếu tố quyết định giúp Dashboard hiển thị thông tin trực quan và chính xác trên bản đồ số.

\subsubsection{Thành viên Đinh Nguyễn Việt Khiêm: Chuyên gia Fact và Tính toàn vẹn dữ liệu}

Với vai trò là chuyên gia phụ trách tầng dữ liệu sự kiện và kiểm soát chất lượng, Đinh Nguyễn Việt Khiêm đảm nhiệm nhiệm vụ then chốt trong việc biến các con số đo lường thô thành các chỉ số có giá trị phân tích và dự báo cao. Khiêm đóng vai trò là "người gác cổng" cho tính chính xác của kho dữ liệu, đảm bảo mọi bản ghi sự kiện đều phản ánh đúng thực tế vận hành của mạng lưới giao thông.

\begin{itemize}
    \item \textbf{Thiết kế và Tối ưu hóa bảng Fact chuyên sâu:} 
    \begin{itemize}[leftmargin=*, nosep]
        \item Tập trung hiện thực và cải tiến bảng \texttt{fact\_traffic\_flow}, xác định độ mịn dữ liệu (Granularity) ở mức 1 phút/bản ghi. Đây là cấu trúc dữ liệu tối ưu giúp cung cấp tập dữ liệu đặc trưng (Features) phong phú cho việc huấn luyện các mô hình AI/ML về sau.
        \item Cải tiến logic của các bảng Fact dựa trên việc phân tích sâu các yêu cầu nghiệp vụ thực tế, đảm bảo các chỉ số về vận tốc, lưu lượng và mật độ dòng xe được tính toán và lưu trữ một cách nhất quán.
        \item Nghiên cứu các xu hướng thiết kế hiện đại để áp dụng vào việc tối ưu hóa Galaxy Schema, giúp các bảng Fact có khả năng liên kết linh hoạt với nhiều nhóm chiều khác nhau.
    \end{itemize}

    \item \textbf{Kiểm soát chất lượng và Đảm bảo tính toàn vẹn (QA \& Integrity):} 
    \begin{itemize}[leftmargin=*, nosep]
        \item Thực hiện các quy trình kiểm thử nghiêm ngặt về tính khả thuộc tính của Schema, đảm bảo mọi khóa ngoại trong bảng Fact đều tham chiếu chính xác đến các bản ghi tương ứng trong bảng Dimension.
        \item Triển khai các cơ chế kiểm tra để loại bỏ triệt để tình trạng "dữ liệu mồ côi" và các sai sót trong báo cáo thống kê PCU (Passenger Car Unit), đảm bảo tính tin cậy tuyệt đối cho các phân tích vĩ mô.
        \item Kiểm tra độ mịn của dữ liệu để đảm bảo không xảy ra hiện tượng mất mát thông tin hoặc trùng lặp bản ghi trong quá trình hợp nhất dữ liệu đa nguồn.
    \end{itemize}

    \item \textbf{Tìm kiếm và Hợp nhất nguồn dữ liệu bổ trợ:} 
    \begin{itemize}[leftmargin=*, nosep]
        \item Chủ động nghiên cứu và tìm kiếm các nguồn dữ liệu từ các nhóm nghiên cứu liên quan và các tệp dữ liệu lưu trữ để làm phong phú thêm kho dữ liệu.
        \item Phân tích tính tương thích của các nguồn dữ liệu mới với cấu trúc Schema hiện tại, đề xuất các phương án tích hợp hiệu quả để mở rộng khả năng phân tích hành vi người dùng và hiệu quả giao thông.
    \end{itemize}
\end{itemize}



\paragraph{Ý nghĩa đóng góp:} Những nỗ lực của Khiêm trong việc kiểm soát chất lượng và thiết kế bảng Fact đã biến kho dữ liệu thành một "mỏ dữ liệu sạch", đạt tiêu chuẩn để đưa vào các mô hình phân tích dự báo. Sự tỉ mỉ trong khâu kiểm tra độ mịn và tính toàn vẹn của Khiêm chính là yếu tố bảo chứng cho sự chính xác của các báo cáo hỗ trợ ra quyết định mà hệ thống cung cấp.

\subsection{Kết chương }

Chương 2 đã cung cấp cái nhìn tổng quan về bối cảnh, động lực và phương pháp luận tiếp cận của đề tài nghiên cứu. Thông qua việc xác định rõ tính cấp thiết của việc xây dựng một hệ thống Kho dữ liệu giao thông thông minh, chương này đã đặt nền móng quan trọng về mặt định hướng triển khai cho toàn bộ đồ án.

Các nội dung trọng tâm đã được trình bày bao gồm:
\textbf{Quản trị và Phân công:} Thiết lập một quy trình làm việc chuyên nghiệp dựa trên các công cụ hiện đại (Trello, GitHub) và phân chia nhiệm vụ chi tiết dựa trên thế mạnh của từng thành viên. Sự phối hợp chặt chẽ giữa các vai trò Kiến trúc sư ETL (Dũng), Chuyên gia GIS (Liêm) và Chuyên gia QA/Fact (Khiêm) chính là chìa khóa đảm bảo tính khả thi cho các giai đoạn tiếp theo.

Với những mục tiêu và kế hoạch đã được xác lập cụ thể, chương tiếp theo sẽ đi sâu vào việc trình bày các \textbf{Kiến thức và Công nghệ nền tảng}, tạo cơ sở lý luận và kỹ thuật vững chắc để hiện thực hóa mô hình Galaxy Schema và Pipeline dữ liệu đã đề ra.